% P5-owned report sections: combined results, ablation discussion, conclusion (draft 9 Nov 2026).
% Budget: about 0.75 + 0.3 + 0.2 page. \todo{[row, column]} cells come from `python -m p5.table` (ablation.tex)
% and from the row JSONs; sentences in \todo{} that depend on a result are chosen once the numbers exist.

\section{Combined Results}
\label{sec:results}

\paragraph{Setup.}
Every row starts from the same frozen artifacts: SimLM passage and query embeddings (passages encoded with
their titles), one BatchTopK SAE with $k=32$ trained with the released objective, and one CL-SR index that
keeps at most 24 concepts per passage (23.61 on average; the paper's CL-SR Efficient reports 22.95). Scoring
uses the released code's BM25 ($k_1=0.6$, $b=2.0$, $k_2=1.5$, BM25 IDF, length as the number of indexed
concepts). Our reproduced baseline B0 reaches 0.316 MRR@10 and 0.947 Recall@1000 on
MS MARCO dev, against 0.343 and 0.954 in the paper; all comparisons are against our B0, not the published
figure. An independent Lucene implementation of the same index (P4) gives 0.3347 MRR@10 on the validation
split against our 0.3336; the small gap comes from its integer impacts and from dropping concepts with
negative IDF, which the released scorer keeps. Each piece was tuned on the 5{,}000
validation queries only; dev and TREC DL 2019/2020 were run once.

\paragraph{Reading the table.}
Table~\ref{tab:ablation} lists B0, one row per piece and \todo{the combined row}. A dagger marks a
difference from B0 with $p<0.05$ (paired randomisation test, 10{,}000 sign flips over all queries, uncorrected).
Every query counts, including those for which CL-SR returns nothing. The cost columns are on the right: query
concepts, index size and latency. Latency is single-threaded wall time of a vectorised re-implementation of the
released retrieval loop that returns identical rankings; the released Python loop takes seconds per query on
the full index. Postings touched per query, which does not depend on the implementation, is reported in the
text.

\paragraph{Which pieces help.}
\todo{Rows in order of [row, MRR@10] gain over B0. For each: the gain, whether it is significant, what it costs
(Concepts/q, Index GB, ms/q). P2, P3 and P4 change passages and move Index GB; P1 and P5 change only scoring or
queries and leave the index unchanged.}

\paragraph{Vocabulary-mismatch queries.}
We rebuild the paper's mismatch sets with Pyserini BM25 ($k_1=0.9$, $b=0.4$) over the prebuilt MS MARCO index:
a query belongs to the top-$K$ set when BM25 retrieves none of its relevant passages in its top $K$. On dev
this gives 4{,}289, 2{,}327 and 988 queries at $K=10$, 100 and 1000 (paper: about 4.3k, 2.3k and 988), and
3{,}196, 1{,}723 and 684 of the 5{,}000 validation queries, which lets us tune on mismatch queries without
touching dev. On the 988-query set B0 reaches 0.104 MRR@10 against SimLM's 0.185 and the paper's 0.124 for
CL-SR Efficient (Table~5). Query expansion, the piece designed for this column, closes little of the gap: the
learned expander reaches 0.107 and a lower query threshold 0.110, neither significant after correction
(Section~\ref{sec:p5-ablation}). \todo{Other pieces' mismatch numbers.}

\paragraph{Do gains add up?}
\todo{Combined row against the sum of single-piece gains, or: the combined run was not completed before the
deadline and the table reports single pieces only.}

\paragraph{Negative results.}
Learned query expansion does not beat B0 after correction and is no better than lowering the query threshold;
it is trained only to reorder SimLM's top 100, so it never learns to recover passages outside them.
Decoder-neighbour expansion lowers every metric, and more neighbours lower it further (validation MRR@10
0.285, 0.259 and 0.227 for budgets 4, 8 and 16), which suggests that neighbours in decoder space mostly match
unrelated passages. \todo{One sentence per other row that does not beat B0, with its reason.}

\subsection{Ablation of Query Expansion}
\label{sec:p5-ablation}

P5's claims rest on six tests, the learned expander L against B0, the query-threshold control T and the
decoder-neighbour control N, each on Recall@1000 and on mismatch@1000 MRR@10, corrected together with Holm's
method at $\alpha=0.05$. T's scale and N's budget are the ones with the highest validation mismatch@1000
MRR@10, so each control is as strong as possible on the claim metric (T at scale 0.5, N at budget 4).
On dev the learned expander does not pass its own test. Against B0 it moves Recall@1000 from 0.947 to 0.949
($p_\text{Holm}=0.24$) and mismatch@1000 MRR@10 from 0.104 to 0.107 ($p_\text{Holm}=0.40$). The threshold
control is at least as good on both and significantly better on Recall@1000 (0.952 against 0.949,
$p_\text{Holm}=0.004$). L beats the neighbour rule on both metrics ($p_\text{Holm}\le 0.006$), but N is itself
below B0 on every metric. On validation, where T and N were chosen, L's mismatch gain over B0 was significant
(0.111 to 0.119, $p_\text{Holm}=0.04$); on dev it is less than half as large. The gain from expansion comes from
near-miss concepts, and a lower query threshold, which needs no training, captures it at least as well; decoder
neighbours mostly add noise. The learned expander agrees: on validation 81\% of the concepts it adds are
near-miss concepts only, 13\% are both near-miss and decoder neighbour, and 6\% are neighbours only (82\%, 12\%
and 6\% of the added weight), so L acts mostly as a learned per-query threshold. It does not raise recall
because it reorders passages the query already reaches: of the 184 validation queries whose relevant passage
B0 misses at 1000, L recovers 11 and loses 11 others. Every one of those 184 relevant passages already shares
an indexed concept with the original query, and 106 contain a concept L adds; they are missed because they
score too low, not because they share nothing with the query. L's training signal points the same way: it
learns to reorder SimLM's top 100, so a relevant passage outside them never affects it.
L adds 16.0 concepts per query on average (39.9 against 23.9) and touches 1.11 times the postings of B0
(442 against 385\,ms per query). T adds 34.8 concepts but touches about as many postings (1.11 times, 446\,ms),
because near-miss concepts are rare and their posting lists short, so L has no cost advantage either. On TREC
DL neither expansion changes NDCG@10 (2019: 0.611 for B0, L and T; 2020: 0.568, 0.574 and 0.569, all
$p>0.3$, uncorrected), while N lowers it on both years (0.522 and 0.512, $p<0.01$).
Across budgets $B\in\{4,8,16\}$ on validation, \todo{recall/latency trade-off}. The tests above use the first training
seed; over three seeds L reaches $0.320\pm0.000$ MRR@10, $0.949\pm0.000$ Recall@1000 and $0.106\pm0.001$
mismatch@1000 MRR@10 on dev, and $0.612\pm0.004$ and $0.572\pm0.003$ NDCG@10 on DL 2019 and 2020 (mean
$\pm$ sd). The seed spread is far smaller than the gap to T, which beats every seed on dev Recall@1000 and
mismatch MRR@10, so the outcome does not depend on the seed.
On the validation mismatch subset, L moves the first relevant passage up for 268 queries and down for 174
(\texttt{p5.example}). The largest gain is ``the definition of protection'', whose relevant dictionary entry
moves from rank 5 to rank 1. The five strongest added concepts are all near-miss concepts, and the training
queries that activate them most suggest they encode the form of a question rather than its topic: the word
``can'', ``what is''/``define'' questions, lists, ``where does \ldots\ come from'' questions, and
``challenges''/``problems''. Only one of the five occurs in the relevant passage. As in the aggregate figures,
the additions re-weight passages the query already matches instead of supplying a missing synonym.

\section{Conclusion}
\label{sec:conclusion}

We asked whether training a stage of concept-level sparse retrieval to rank, instead of leaving it hand-set,
narrows the gap to the dense retriever it is built from. \todo{Answer in numbers: best [*, MRR@10] and
[*, R@1k] against B0 and against SimLM's 0.411 (paper Table~1), and what it costs in index size and latency.}
\todo{One sentence on which kind of change helps most (scoring, SAE training, sparsity, pruning or query
expansion).}

\paragraph{Limitations.}
We study one SAE size ($k=32$) on MS MARCO only, without BEIR transfer. Latency is measured single-threaded
on shared NSCC nodes. Final trained rows use three seeds. The mismatch sets come from our Pyserini run, which
matches the paper's dev count at top 1000 exactly but may differ in individual queries. The best threshold scale
(0.5) and the learned expander's budget (16) are both at the edge of the ranges we searched.
